% GENERATED FILE. Edit canonical lessons, then run npm run generate:books.
% canonical-source-hash: fnv1a64:930fb7c410b7837c
% canonical-lessons: GU-C08-vicharvun, GU-C08-samajvun, GU-C08-vanchvun, GU-C08-lakhvun

\chapter{The Mind and the Page}
\label{ch:gu-mind-and-page}
\hlchaptermodality{10}

\begin{chapteropening}
\textbf{By the end of this chapter:} \emph{I can say that I think, understand, read and write in Gujarati, hear and write the nasal that sits in the middle of \gu{વાંચવું}, and name what each of the four roots was doing before it meant a mental act — ranging about, waking up, speaking, and scratching.}

Say \gu{હું} \gu{મારું} \gu{નામ} \gu{લખું} \gu{છું}, run the chapter's four verbs back with the four things their roots meant — vichārvũ from a root meaning to range about, samajvũ from one meaning to wake, vā̃chvũ from one meaning to speak, lakhvũ from one meaning to scratch — place the medial nasal of vā̃ch- and the assimilate-then-simplify path from dhy to j, and say which of the four honestly has no English cousin at all.
\end{chapteropening}

\section[vichārvũ]{\gu{વિચારવું} — \textquotedblleft{}to think,\textquotedblright{} which began as wandering}

\label{lesson:GU-C08-vicharvun}



\begin{quote}
\emph{Jāṇvũ} — take the ending off it, and name the plain English verb that is the same word. Knowing is where you arrive. Here is what you do on the way.
\end{quote}

\subsection*{You'll want to know first — one word}
\begin{quote}
\textbf{\gu{વિચારવું}} — \emph{vichārvũ} — \textbf{to think}
\end{quote}

The stem is \textbf{\gu{વિચાર}-} \emph{vichār-}: \emph{hũ vichārũ chhũ}, \textquotedblleft{}I think.\textquotedblright{}

\subsection*{The letters in this word}
\textbf{\gu{વિ}} is \textbf{\gu{વ}} \emph{va} wearing \textbf{\gu{િ}}, the short-\emph{i} sign — and that sign stands to the \textbf{left} of its letter although the sound comes after it. Then \emph{chā}: the letter \textbf{\gu{ચ}} \emph{cha}, already yours from \emph{chār}, \textquotedblleft{}four,\textquotedblright{} under the long-\emph{ā} sign. Then \textbf{\gu{ર}} \emph{ra}, and the ending \textbf{\gu{વું}}.

\begin{cousinweb}
\textbf{\gu{વિચાર}} \emph{vichār} is a word before it is a verb: it means \textbf{a thought}. It is Sanskrit \emph{vicāra} — \textbf{vi-} \textquotedblleft{}apart, thoroughly\textquotedblright{} on the root \textbf{car-}, \textquotedblleft{}to move, to range about.\textquotedblright{} A thought is named as a \textbf{ranging over}, the mind walking round a thing.

The root is Indo-European, *\emph{kwel-} \textquotedblleft{}to turn.\textquotedblright{} Germanic wore it into \textbf{wheel}, whose whole business is turning. Greek made \emph{kúklos}, \textquotedblleft{}a circle,\textquotedblright{} which English borrowed as \textbf{cycle} and \textbf{encyclopedia}. Latin made \emph{colere}, \textquotedblleft{}to till, to dwell in,\textquotedblright{} behind \textbf{cultivate}, \textbf{culture} and \textbf{colony}.
\end{cousinweb}

\begin{grammarlens}[title={What you've built: a verb Gujarati made rather than inherited}]
\emph{Hovũ}, \emph{javũ}, \emph{āvvũ}, \emph{khāvũ}, \emph{jovũ} and \emph{jāṇvũ} came down the ordinary road, worn thinner at every step. \emph{Vichārvũ} did not: Gujarati took the \textbf{noun} \emph{vichār} and put a verb ending on it, on the model of Sanskrit \emph{vicarati}.

You have watched Sanskrit reach forward like this before. Count to three: \emph{ek, be, traṇ}. Prakrit had lost the \emph{r} of \emph{trīṇi}, and \emph{traṇ} carries it because Sanskrit put it back. A language that never stops reading its ancestor never stops borrowing from it. \emph{Traṇ} got a sound that way; \emph{vichārvũ} got a whole verb.
\end{grammarlens}

\subsection*{Your turn}
\begin{itemize}
\raggedright
  \item \emph{Say it:} \textbf{vichārvũ} — to think
  \item \emph{Strip:} minus \textbf{-vũ} leaves \textbf{vichār-}, which is also the noun \textquotedblleft{}a thought\textquotedblright{}
  \item \emph{Build:} \textbf{hũ vichārũ chhũ} — I think
  \item \emph{Trace it:} \emph{vi-} plus \emph{car-} \textquotedblleft{}range about\textquotedblright{} $\to$ \textbf{wheel}, \textbf{cycle}, \textbf{colony}
  \item \emph{Contrast:} \textbf{jāṇvũ} handed down, and English \textbf{know}; \textbf{vichārvũ} built — as \emph{traṇ} had its \emph{r} handed back
\end{itemize}

\begin{tcolorbox}[breakable,colback=teal!4,colframe=teal!35!black,title={Before you move on}]
What did the root behind \emph{vichār} mean first, and which everyday English word turns on it? (\textquotedblleft{}To range about\textquotedblright{}; \textbf{wheel}.) Was \emph{vichārvũ} handed down or built, and which count of yours was patched the same way? (Built on the noun \emph{vichār}; \emph{traṇ} had its \emph{r} restored.) Say \textquotedblleft{}I think.\textquotedblright{} (\emph{Hũ vichārũ chhũ}.)

Source: \href{https://en.wiktionary.org/wiki/\%E0\%AA\%B5\%E0\%AA\%BF\%E0\%AA\%9A\%E0\%AA\%BE\%E0\%AA\%B0\%E0\%AA\%B5\%E0\%AB\%81\%E0\%AA\%82}{Wiktionary: \gu{વિચારવું}}.
\end{tcolorbox}
\section[samajvũ]{\gu{સમજવું} — \textquotedblleft{}to understand,\textquotedblright{} which is waking up}

\label{lesson:GU-C08-samajvun}



\begin{quote}
Two verbs and one root meaning: \emph{jāṇvũ}, and \emph{vichārvũ} on a root that meant \textbf{to range about}. The third sits between them, with a third kind of picture again.
\end{quote}

\subsection*{You'll want to know first — one word}
\begin{quote}
\textbf{\gu{સમજવું}} — \emph{samajvũ} — \textbf{to understand}
\end{quote}

The stem is \textbf{\gu{સમજ}-} \emph{samaj-}: \emph{hũ samajũ chhũ}, \textquotedblleft{}I understand.\textquotedblright{} \emph{Samaj} on its own, said alone, means \textquotedblleft{}understood.\textquotedblright{}

\subsection*{The letters in this word}
Not one new letter. \textbf{\gu{સ}} \emph{sa} stood in \emph{namaste}, \textbf{\gu{મ}} \emph{ma} in every \emph{mārũ}, \textbf{\gu{જ}} \emph{ja} opened \emph{javũ}. Three bare consonants carrying their built-in \emph{a}, then the ending: \emph{sa-ma-ja-vũ}.

\begin{cousinweb}
\textbf{samajvũ} comes from Sanskrit \emph{sambudhyate}: \textbf{sam-} \textquotedblleft{}together, fully,\textquotedblright{} on the root \textbf{budh-}, \textquotedblleft{}\textbf{to wake, to become aware}.\textquotedblright{} To understand is to come awake to a thing.

That root named a person too. One who has fully woken is a \textbf{buddha} — the title is this root with a participle ending, and it means exactly \textquotedblleft{}the awakened one.\textquotedblright{}

It is Indo-European as well, *\emph{bheudh-}, and English keeps it in quieter clothes. Old English \emph{bēodan}, \textquotedblleft{}to announce, to offer,\textquotedblright{} left \textbf{bode} and \textbf{forebode} — to bode ill is to give warning, which is what a waking mind does — and \textbf{forbid}, an announcement pointed the other way.
\end{cousinweb}

\begin{grammarlens}[title={What you've built: a cluster that assimilated, then simplified}]
\emph{Sambudhyate} has a \textbf{dhy} in the middle; Gujarati has a plain \textbf{j}. Two steps did that, in this order. \textbf{Assimilate:} the \emph{dh} was remade in the shape of the \emph{y} after it, both landing on the palate, and the pair came out doubled as \textbf{jjh} — Prakrit \emph{saṃbujjhadi}. \textbf{Simplify:} the double then fell to a single \textbf{j}, giving \emph{samaj-}.

You have watched this pair once already. Count to two: \emph{be}. Its \emph{dv-} went to \textbf{bb} because the \emph{d} was remade in the shape of the \emph{v} beside it, and the double then simplified to \textbf{b}. Different letters, same two moves.
\end{grammarlens}

\subsection*{Your turn}
\begin{itemize}
\raggedright
  \item \emph{Say it:} \textbf{samajvũ}, then \textbf{hũ samajũ chhũ} — I understand
  \item \emph{Trace it:} \emph{sambudhyate} $\to$ \emph{saṃbujjhadi} $\to$ \textbf{samaj-}: assimilate, then simplify
  \item \emph{Match:} the same two moves in \emph{dv} $\to$ \emph{bb} $\to$ \textbf{b}, which gave you \emph{be}
  \item \emph{Connect:} \textbf{budh-} \textquotedblleft{}to wake\textquotedblright{} $\to$ \textbf{buddha}, and English \textbf{bode}, \textbf{forbid}
  \item \emph{Run through:} \textbf{jāṇvũ}, \textbf{vichārvũ}, \textbf{samajvũ} — know, think, understand
\end{itemize}

\begin{tcolorbox}[breakable,colback=teal!4,colframe=teal!35!black,title={Before you move on}]
What did the root behind \emph{samajvũ} mean, and which famous title is the same root? (\textquotedblleft{}To wake, to become aware\textquotedblright{}; \textbf{Buddha}.) Name the two steps from \emph{dhy} to \emph{j}, and the number that took the same two. (Assimilate, then simplify; \emph{be}, from \emph{dv} to \emph{bb} to \emph{b}.) Say \textquotedblleft{}I understand.\textquotedblright{} (\emph{Hũ samajũ chhũ}.)

Source: \href{https://en.wiktionary.org/wiki/\%E0\%AA\%B8\%E0\%AA\%AE\%E0\%AA\%9C\%E0\%AA\%B5\%E0\%AB\%81\%E0\%AA\%82}{Wiktionary: \gu{સમજવું}}.
\end{tcolorbox}
\section[vā̃chvũ]{\gu{વાંચવું} — \textquotedblleft{}to read,\textquotedblright{} which is making a page speak}

\label{lesson:GU-C08-vanchvun}



\begin{quote}
\emph{Hũ samajũ chhũ} — and what did that root mean first? Then \emph{vichārvũ}. The next verb puts the words in front of you.
\end{quote}

\subsection*{You'll want to know first — one word}
\begin{quote}
\textbf{\gu{વાંચવું}} — \emph{vā̃chvũ} — \textbf{to read}
\end{quote}

The stem is \textbf{\gu{વાંચ}-} \emph{vā̃ch-}, and nothing new happens after it: \emph{hũ vā̃chũ chhũ}, \textquotedblleft{}I read.\textquotedblright{}

\subsection*{The letters in this word}
\textbf{\gu{વા}} is \textbf{\gu{વ}} \emph{va} with the long-\emph{ā} sign, and the dot \textbf{\gu{ં}} sits on top of it. Then \textbf{\gu{ચ}} \emph{cha}, the letter you counted \emph{chār} with and thought with in \emph{vichār}. Then \textbf{\gu{વું}}. \emph{Vā̃-cha-vũ} — and the dot in the middle is the piece to be careful about.

\begin{grammarlens}[title={What you've built: a nasal that is not at the end}]
Every nasal dot so far has been at the \textbf{end} of a word: \emph{hũ}, \emph{chhũ}, and the \textbf{-\gu{વું}} on every verb. This one is in the \textbf{middle}, doing the same job in a new place — the \emph{ā} before \textbf{\gu{ચ}} goes through the nose, and it must be heard as well as written. Count to five and stop on \textbf{\gu{પાંચ}} \emph{pā̃ch}: same dot, same position, same nasal \emph{ā} running into a \textbf{\gu{ચ}}.

And you know how much one small mark can carry. \textbf{\gu{જવું}} \emph{javũ} is \textquotedblleft{}to go\textquotedblright{} and \textbf{\gu{જોવું}} \emph{jovũ} is \textquotedblleft{}to see,\textquotedblright{} and the whole difference is the sign on the \textbf{\gu{જ}}. This dot is the same kind of hinge: small enough to skip, and not skippable.
\end{grammarlens}

\begin{cousinweb}
\textbf{vā̃chvũ} comes through Old Gujarati from Sanskrit \emph{vācayati}, a \textbf{causative}: the root \textbf{vac-} means \textquotedblleft{}to speak,\textquotedblright{} so \emph{vācayati} is \textquotedblleft{}\textbf{to make speak}.\textquotedblright{} Reading was named as making a silent page talk — which is what reading was, aloud, for most of the time people have done it.

The root is Indo-European, *\emph{wekw-}, and Latin holds it as \emph{vōx}, \textquotedblleft{}voice.\textquotedblright{} English took the family whole: \textbf{voice}, \textbf{vocal}, \textbf{vowel}, \textbf{vocation} (a calling), \textbf{advocate} (one called to your side), \textbf{provoke}, \textbf{revoke}.
\end{cousinweb}

\subsection*{Your turn}
\begin{itemize}
\raggedright
  \item \emph{Say it:} \textbf{vā̃chvũ}, then \textbf{hũ vā̃chũ chhũ} — I read
  \item \emph{Match:} the same middle nasal in \textbf{pā̃ch}, \textquotedblleft{}five\textquotedblright{}
  \item \emph{Contrast:} \textbf{javũ} and \textbf{jovũ} — one mark decides which word it is
  \item \emph{Trace it:} \emph{vac-} \textquotedblleft{}speak\textquotedblright{} $\to$ \emph{vācayati} \textquotedblleft{}make speak\textquotedblright{} $\to$ \textbf{vā̃chvũ}, and $\to$ \textbf{voice}, \textbf{vowel}, \textbf{advocate}
  \item \emph{Run through:} \textbf{vichārvũ}, \textbf{samajvũ}, \textbf{vā̃chvũ}
\end{itemize}

\begin{tcolorbox}[breakable,colback=teal!4,colframe=teal!35!black,title={Before you move on}]
Where does the nasal dot sit in \emph{vā̃chvũ}, and which number of yours carries it in the same place? (In the middle, on the \emph{ā}; \emph{pā̃ch}.) What did \emph{vācayati} mean literally, and which English word for the thing you speak with is its cousin? (\textquotedblleft{}To make speak\textquotedblright{}; \textbf{voice}.) Say \textquotedblleft{}I read.\textquotedblright{} (\emph{Hũ vā̃chũ chhũ}.)

Source: \href{https://en.wiktionary.org/wiki/\%E0\%AA\%B5\%E0\%AA\%BE\%E0\%AA\%82\%E0\%AA\%9A\%E0\%AA\%B5\%E0\%AB\%81\%E0\%AA\%82}{Wiktionary: \gu{વાંચવું}}.
\end{tcolorbox}
\section[lakhvũ]{\gu{લખવું} — \textquotedblleft{}to write,\textquotedblright{} which began as scratching}

\label{lesson:GU-C08-lakhvun}



\begin{quote}
Three verbs, three pictures: \emph{vichārvũ} on a root meaning \textbf{to range about}, \emph{samajvũ} on one meaning \textbf{to wake}, \emph{vā̃chvũ} on one meaning \textbf{to speak}. The fourth root is the most physical of them.
\end{quote}

\subsection*{You'll want to know first — one word}
\begin{quote}
\textbf{\gu{લખવું}} — \emph{lakhvũ} — \textbf{to write}
\end{quote}

The stem is \textbf{\gu{લખ}-} \emph{lakh-}, and now you can say something worth saying:

\textbf{\gu{હું} \gu{મારું} \gu{નામ} \gu{લખું} \gu{છું}} — \emph{hũ mārũ nām lakhũ chhũ} — \textquotedblleft{}\textbf{I write my name}.\textquotedblright{}

\subsection*{The letters in this word}
Both letters are already yours. \textbf{\gu{લ}} \emph{la} stood in \emph{bolvũ}, \textquotedblleft{}to speak.\textquotedblright{} \textbf{\gu{ખ}} \emph{kha} is the puffed \emph{k} of \emph{khāvũ}, \textquotedblleft{}to eat.\textquotedblright{} \emph{La-kha-vũ}.

\begin{cousinweb}
\textbf{lakhvũ} comes through Prakrit from Sanskrit \emph{likhati}, and \textbf{likh-} meant \textquotedblleft{}\textbf{to scratch, to score}\textquotedblright{} long before it meant to write. Its plainest Sanskrit relative says so out loud: \textbf{\gu{રેખા}} \emph{rekhā}, \textquotedblleft{}a line, a streak.\textquotedblright{}

Beyond Indo-Aryan the cousins are disputed, and this lesson claims none — the same answer \emph{khāvũ} gave. What is not disputed is that three families reached for one picture: Latin \emph{scrībere} also meant \textbf{to scratch}, behind \textbf{scribe} and \textbf{script}, and English \textbf{write} began as Germanic \textquotedblleft{}to score.\textquotedblright{}
\end{cousinweb}

\begin{grammarlens}[title={What you've built: four verbs, and what each root did first}]
\noindent
\begin{tabularx}{\linewidth}{@{}>{\raggedright\arraybackslash}X>{\raggedright\arraybackslash}X@{}}
\toprule
\textbf{verb} & \textbf{what its root did first} \\
\midrule
\emph{vichārvũ} — think & ranged about, turned \\
\emph{samajvũ} — understand & woke up \\
\emph{vā̃chvũ} — read & spoke \\
\emph{lakhvũ} — write & scratched \\
\bottomrule
\end{tabularx}

Not one of the four began as a mental act. The mind gets its vocabulary from the body, every time.

Three habits ran through the chapter. A cluster \textbf{assimilates and then simplifies} — \emph{dhy} to \emph{jjh} to \emph{j} in \emph{samajvũ}, whose root \emph{budh-} also gave \textbf{buddha}. A \textbf{dot in the middle} is a nasal you must hear — \emph{vā̃ch-}, whose root \emph{vac-} also gave \textbf{voice}. And when the trail stops you say so: \emph{lakh-} and \emph{khāvũ} have no secure English cousin, \emph{jovũ}'s origin is probable rather than proven, and \emph{jāṇvũ} is the one that paid out, being English \textbf{know}.
\end{grammarlens}

\subsection*{Your turn}
\begin{itemize}
\raggedright
  \item \emph{Produce:} \textbf{hũ mārũ nām lakhũ chhũ} — I write my name
  \item \emph{Run through:} \textbf{vichārvũ}, \textbf{samajvũ}, \textbf{vā̃chvũ}, \textbf{lakhvũ} — stripping \textbf{-vũ} off each
  \item \emph{Name:} what each root did first — ranged, woke, spoke, scratched
  \item \emph{Say it:} the two steps behind \emph{samaj-}, and where the dot sits in \textbf{vā̃chvũ}
  \item \emph{Connect:} \textbf{rekhā}, still meaning the scratch
\end{itemize}

\begin{tcolorbox}[breakable,colback=teal!4,colframe=teal!35!black,title={Before you move on}]
Say \textquotedblleft{}I write my name.\textquotedblright{} (\emph{Hũ mārũ nām lakhũ chhũ}.) Name the four verbs and what each root did before it meant a mental act. (Ranged, woke, spoke, scratched.) Which Sanskrit word is \emph{lakh-}'s plainest relative? (\textbf{Rekhā}, \textquotedblleft{}a line.\textquotedblright{}) Which of your verbs has a certain English cousin, and which have none? (\emph{Jāṇvũ} is \textbf{know}; \emph{lakh-} and \emph{khāvũ} have none.)

Source: \href{https://en.wiktionary.org/wiki/\%E0\%AA\%B2\%E0\%AA\%96\%E0\%AA\%B5\%E0\%AB\%81\%E0\%AA\%82}{Wiktionary: \gu{લખવું}}.
\end{tcolorbox}
